\IfFileExists{stacks-project.cls}{%
\documentclass{stacks-project}
}{%
\documentclass{amsart}
}

% For dealing with references we use the comment environment
\usepackage{verbatim}
\newenvironment{reference}{\comment}{\endcomment}
%\newenvironment{reference}{}{}
\newenvironment{slogan}{\comment}{\endcomment}
\newenvironment{history}{\comment}{\endcomment}

% For commutative diagrams we use Xy-pic
\usepackage[all]{xy}

% We use 2cell for 2-commutative diagrams.
\xyoption{2cell}
\UseAllTwocells

% We use multicol for the list of chapters between chapters
\usepackage{multicol}

% This is generally recommended for better output
\usepackage{lmodern}
\usepackage[T1]{fontenc}

% For cross-file-references
\usepackage{xr-hyper}

% Package for hypertext links:
\usepackage{hyperref}

% For any local file, say "hello.tex" you want to link to please
% use \externaldocument[hello-]{hello}
\externaldocument[introduction-]{introduction}
\externaldocument[conventions-]{conventions}
\externaldocument[sets-]{sets}
\externaldocument[categories-]{categories}
\externaldocument[topology-]{topology}
\externaldocument[sheaves-]{sheaves}
\externaldocument[sites-]{sites}
\externaldocument[stacks-]{stacks}
\externaldocument[fields-]{fields}
\externaldocument[algebra-]{algebra}
\externaldocument[brauer-]{brauer}
\externaldocument[homology-]{homology}
\externaldocument[derived-]{derived}
\externaldocument[simplicial-]{simplicial}
\externaldocument[more-algebra-]{more-algebra}
\externaldocument[smoothing-]{smoothing}
\externaldocument[modules-]{modules}
\externaldocument[sites-modules-]{sites-modules}
\externaldocument[injectives-]{injectives}
\externaldocument[cohomology-]{cohomology}
\externaldocument[sites-cohomology-]{sites-cohomology}
\externaldocument[dga-]{dga}
\externaldocument[dpa-]{dpa}
\externaldocument[sdga-]{sdga}
\externaldocument[hypercovering-]{hypercovering}
\externaldocument[schemes-]{schemes}
\externaldocument[constructions-]{constructions}
\externaldocument[properties-]{properties}
\externaldocument[morphisms-]{morphisms}
\externaldocument[coherent-]{coherent}
\externaldocument[divisors-]{divisors}
\externaldocument[limits-]{limits}
\externaldocument[varieties-]{varieties}
\externaldocument[topologies-]{topologies}
\externaldocument[descent-]{descent}
\externaldocument[perfect-]{perfect}
\externaldocument[more-morphisms-]{more-morphisms}
\externaldocument[flat-]{flat}
\externaldocument[groupoids-]{groupoids}
\externaldocument[more-groupoids-]{more-groupoids}
\externaldocument[etale-]{etale}
\externaldocument[chow-]{chow}
\externaldocument[intersection-]{intersection}
\externaldocument[pic-]{pic}
\externaldocument[weil-]{weil}
\externaldocument[adequate-]{adequate}
\externaldocument[dualizing-]{dualizing}
\externaldocument[duality-]{duality}
\externaldocument[discriminant-]{discriminant}
\externaldocument[derham-]{derham}
\externaldocument[local-cohomology-]{local-cohomology}
\externaldocument[algebraization-]{algebraization}
\externaldocument[curves-]{curves}
\externaldocument[resolve-]{resolve}
\externaldocument[models-]{models}
\externaldocument[functors-]{functors}
\externaldocument[equiv-]{equiv}
\externaldocument[pione-]{pione}
\externaldocument[etale-cohomology-]{etale-cohomology}
\externaldocument[proetale-]{proetale}
\externaldocument[relative-cycles-]{relative-cycles}
\externaldocument[more-etale-]{more-etale}
\externaldocument[trace-]{trace}
\externaldocument[crystalline-]{crystalline}
\externaldocument[spaces-]{spaces}
\externaldocument[spaces-properties-]{spaces-properties}
\externaldocument[spaces-morphisms-]{spaces-morphisms}
\externaldocument[decent-spaces-]{decent-spaces}
\externaldocument[spaces-cohomology-]{spaces-cohomology}
\externaldocument[spaces-limits-]{spaces-limits}
\externaldocument[spaces-divisors-]{spaces-divisors}
\externaldocument[spaces-over-fields-]{spaces-over-fields}
\externaldocument[spaces-topologies-]{spaces-topologies}
\externaldocument[spaces-descent-]{spaces-descent}
\externaldocument[spaces-perfect-]{spaces-perfect}
\externaldocument[spaces-more-morphisms-]{spaces-more-morphisms}
\externaldocument[spaces-flat-]{spaces-flat}
\externaldocument[spaces-groupoids-]{spaces-groupoids}
\externaldocument[spaces-more-groupoids-]{spaces-more-groupoids}
\externaldocument[bootstrap-]{bootstrap}
\externaldocument[spaces-pushouts-]{spaces-pushouts}
\externaldocument[spaces-chow-]{spaces-chow}
\externaldocument[groupoids-quotients-]{groupoids-quotients}
\externaldocument[spaces-more-cohomology-]{spaces-more-cohomology}
\externaldocument[spaces-simplicial-]{spaces-simplicial}
\externaldocument[spaces-duality-]{spaces-duality}
\externaldocument[formal-spaces-]{formal-spaces}
\externaldocument[restricted-]{restricted}
\externaldocument[spaces-resolve-]{spaces-resolve}
\externaldocument[formal-defos-]{formal-defos}
\externaldocument[defos-]{defos}
\externaldocument[cotangent-]{cotangent}
\externaldocument[examples-defos-]{examples-defos}
\externaldocument[algebraic-]{algebraic}
\externaldocument[examples-stacks-]{examples-stacks}
\externaldocument[stacks-sheaves-]{stacks-sheaves}
\externaldocument[criteria-]{criteria}
\externaldocument[artin-]{artin}
\externaldocument[quot-]{quot}
\externaldocument[stacks-properties-]{stacks-properties}
\externaldocument[stacks-morphisms-]{stacks-morphisms}
\externaldocument[stacks-limits-]{stacks-limits}
\externaldocument[stacks-cohomology-]{stacks-cohomology}
\externaldocument[stacks-perfect-]{stacks-perfect}
\externaldocument[stacks-introduction-]{stacks-introduction}
\externaldocument[stacks-more-morphisms-]{stacks-more-morphisms}
\externaldocument[stacks-geometry-]{stacks-geometry}
\externaldocument[moduli-]{moduli}
\externaldocument[moduli-curves-]{moduli-curves}
\externaldocument[examples-]{examples}
\externaldocument[exercises-]{exercises}
\externaldocument[guide-]{guide}
\externaldocument[desirables-]{desirables}
\externaldocument[coding-]{coding}
\externaldocument[obsolete-]{obsolete}
\externaldocument[fdl-]{fdl}
\externaldocument[index-]{index}

% Theorem environments.
%
\theoremstyle{plain}
\newtheorem{theorem}[subsection]{Theorem}
\newtheorem{proposition}[subsection]{Proposition}
\newtheorem{lemma}[subsection]{Lemma}

\theoremstyle{definition}
\newtheorem{definition}[subsection]{Definition}
\newtheorem{example}[subsection]{Example}
\newtheorem{exercise}[subsection]{Exercise}
\newtheorem{situation}[subsection]{Situation}

\theoremstyle{remark}
\newtheorem{remark}[subsection]{Remark}
\newtheorem{remarks}[subsection]{Remarks}

\numberwithin{equation}{subsection}

% Macros
%
\def\lim{\mathop{\mathrm{lim}}\nolimits}
\def\colim{\mathop{\mathrm{colim}}\nolimits}
\def\Spec{\mathop{\mathrm{Spec}}}
\def\Hom{\mathop{\mathrm{Hom}}\nolimits}
\def\Ext{\mathop{\mathrm{Ext}}\nolimits}
\def\SheafHom{\mathop{\mathcal{H}\!\mathit{om}}\nolimits}
\def\SheafExt{\mathop{\mathcal{E}\!\mathit{xt}}\nolimits}
\def\Sch{\mathit{Sch}}
\def\Mor{\mathop{\mathrm{Mor}}\nolimits}
\def\Ob{\mathop{\mathrm{Ob}}\nolimits}
\def\Sh{\mathop{\mathit{Sh}}\nolimits}
\def\NL{\mathop{N\!L}\nolimits}
\def\CH{\mathop{\mathrm{CH}}\nolimits}
\def\proetale{{pro\text{-}\acute{e}tale}}
\def\etale{{\acute{e}tale}}
\def\QCoh{\mathit{QCoh}}
\def\Ker{\mathop{\mathrm{Ker}}}
\def\Im{\mathop{\mathrm{Im}}}
\def\Coker{\mathop{\mathrm{Coker}}}
\def\Coim{\mathop{\mathrm{Coim}}}

% Boxtimes
%
\DeclareMathSymbol{\boxtimes}{\mathbin}{AMSa}{"02}

%
% Macros for moduli stacks/spaces
%
\def\QCohstack{\mathcal{QC}\!\mathit{oh}}
\def\Cohstack{\mathcal{C}\!\mathit{oh}}
\def\Spacesstack{\mathcal{S}\!\mathit{paces}}
\def\Quotfunctor{\mathrm{Quot}}
\def\Hilbfunctor{\mathrm{Hilb}}
\def\Curvesstack{\mathcal{C}\!\mathit{urves}}
\def\Polarizedstack{\mathcal{P}\!\mathit{olarized}}
\def\Complexesstack{\mathcal{C}\!\mathit{omplexes}}
% \Pic is the operator that assigns to X its picard group, usage \Pic(X)
% \Picardstack_{X/B} denotes the Picard stack of X over B
% \Picardfunctor_{X/B} denotes the Picard functor of X over B
\def\Pic{\mathop{\mathrm{Pic}}\nolimits}
\def\Picardstack{\mathcal{P}\!\mathit{ic}}
\def\Picardfunctor{\mathrm{Pic}}
\def\Deformationcategory{\mathcal{D}\!\mathit{ef}}
\usepackage{cleveref}
\crefname{originallemma}{original lemma}{original lemmas}
\Crefname{originallemma}{Original lemma}{Original lemmas}
\makeatletter
\newlabel{lemma-deform-koszul}{{0G0N}{source}{}{sourcepointer-1}{}}
\newlabel{lemma-deform-koszul@cref}{{[originallemma][0][]0G0N}{[1][0][]source}{}{}{}}
\makeatother
\crefname{originallemma}{original lemma}{original lemmas}
\Crefname{originallemma}{Original lemma}{Original lemmas}
\makeatletter
\newlabel{lemma-isomorphism-in-neighbourhood}{{0G0Q}{source}{}{sourcepointer-2}{}}
\newlabel{lemma-isomorphism-in-neighbourhood@cref}{{[originallemma][0][]0G0Q}{[1][0][]source}{}{}{}}
\makeatother
\crefname{originallemma}{original lemma}{original lemmas}
\Crefname{originallemma}{Original lemma}{Original lemmas}
\makeatletter
\newlabel{lemma-limit-arguments}{{0G0P}{source}{}{sourcepointer-3}{}}
\newlabel{lemma-limit-arguments@cref}{{[originallemma][0][]0G0P}{[1][0][]source}{}{}{}}
\makeatother
\crefname{originallemma}{original lemma}{original lemmas}
\Crefname{originallemma}{Original lemma}{Original lemmas}
\makeatletter
\newlabel{lemma-orthogonal-point-sheaf}{{0G03}{source}{}{sourcepointer-4}{}}
\newlabel{lemma-orthogonal-point-sheaf@cref}{{[originallemma][0][]0G03}{[1][0][]source}{}{}{}}
\makeatother
\crefname{originalreference}{original reference}{original references}
\Crefname{originalreference}{Original reference}{Original references}
\makeatletter
\newlabel{more-algebra-lemma-regular-koszul-regular}{{062F}{source}{}{sourcepointer-5}{}}
\newlabel{more-algebra-lemma-regular-koszul-regular@cref}{{[originalreference][0][]062F}{[1][0][]source}{}{}{}}
\makeatother
\crefname{originalreference}{original reference}{original references}
\Crefname{originalreference}{Original reference}{Original references}
\makeatletter
\newlabel{more-morphisms-lemma-smooth-diagonal-perfect}{{0FDP}{source}{}{sourcepointer-6}{}}
\newlabel{more-morphisms-lemma-smooth-diagonal-perfect@cref}{{[originalreference][0][]0FDP}{[1][0][]source}{}{}{}}
\makeatother
\title{Rigidity of smooth proper fibres with diagonal kernel restriction}
\author{The Stacks Project Authors}
\date{stacks-project:equiv.tex:local-snapshot-20261012; GNU FDL 1.2 or later}
\begin{document}
\maketitle
\noindent\textit{Editorial note.} Two forward references to same-chapter lemma-limit-arguments (judged citation-only this batch, not carded); transcribed as-is.\par\medskip
\noindent\textit{Allowed standard prerequisites (editorial).}\begin{itemize}
\item Orthogonal point sheaf. lemma-orthogonal-point-sheaf (carded this batch, L26488). Case I: kernel vanishes on an open neighbourhood off the diagonal.
\item Koszul deformation lemma. lemma-deform-koszul (carded this batch, L26463). Case II: lifts the Koszul representation over the special fibre.
\item Limit arguments glue. lemma-limit-arguments (same chapter, not carded). Spreads the local Koszul representation to an affine neighbourhood.
\item Isomorphism in a neighbourhood. lemma-isomorphism-in-neighbourhood (carded this batch, L26474). Final step: fibre isomorphism spreads over the base.
\end{itemize}\medskip
\ifdefined\hypertarget\hypertarget{sourcepointer-1}{}\fi
\noindent\textit{Source reference.} Original lemma 0G0N: sources/stacks-equiv/pinned-original-source.tex, source lines 3962--3962. Navigation pointer only: the counted proof cites this statement by \textbackslash{}ref as a standing input; the deduction the counted proof has to make is carried out inside the counted span, and no part of that deduction is delegated to this pointer.\par\medskip
\ifdefined\hypertarget\hypertarget{sourcepointer-2}{}\fi
\noindent\textit{Source reference.} Original lemma 0G0Q: sources/stacks-equiv/pinned-original-source.tex, source lines 4080--4080. Navigation pointer only: the counted proof cites this statement by \textbackslash{}ref as a standing input; the deduction the counted proof has to make is carried out inside the counted span, and no part of that deduction is delegated to this pointer.\par\medskip
\ifdefined\hypertarget\hypertarget{sourcepointer-3}{}\fi
\noindent\textit{Source reference.} Original lemma 0G0P: sources/stacks-equiv/pinned-original-source.tex, source lines 4039--4039. Navigation pointer only: the counted proof cites this statement by \textbackslash{}ref as a standing input; the deduction the counted proof has to make is carried out inside the counted span, and no part of that deduction is delegated to this pointer.\par\medskip
\ifdefined\hypertarget\hypertarget{sourcepointer-4}{}\fi
\noindent\textit{Source reference.} Original lemma 0G03: sources/stacks-equiv/pinned-original-source.tex, source lines 2795--2795. Navigation pointer only: the counted proof cites this statement by \textbackslash{}ref as a standing input; the deduction the counted proof has to make is carried out inside the counted span, and no part of that deduction is delegated to this pointer.\par\medskip
\ifdefined\hypertarget\hypertarget{sourcepointer-5}{}\fi
\noindent\textit{Source reference.} Original reference 062F: focused-seeds/geometry/S9-eqv-no-deformations/ref-tags.txt, source lines 5--5. Navigation pointer only: the counted proof cites this statement by \textbackslash{}ref as a standing input; the deduction the counted proof has to make is carried out inside the counted span, and no part of that deduction is delegated to this pointer.\par\medskip
\ifdefined\hypertarget\hypertarget{sourcepointer-6}{}\fi
\noindent\textit{Source reference.} Original reference 0FDP: focused-seeds/geometry/S9-eqv-no-deformations/ref-tags.txt, source lines 6--6. Navigation pointer only: the counted proof cites this statement by \textbackslash{}ref as a standing input; the deduction the counted proof has to make is carried out inside the counted span, and no part of that deduction is delegated to this pointer.\par\medskip
\begin{lemma}
\label{lemma-no-deformations}
Let $k$ be a field. Let $S$ be a finite type scheme over $k$
with $k$-rational point $s$. Let $Y \to S$ be a smooth proper morphism.
Let $X = Y_s \times S \to S$ be the constant family with fibre
$Y_s$. Let $K$ be the Fourier-Mukai kernel of a relative equivalence
from $X$ to $Y$ over $S$. Assume the restriction
$$
L(Y_s \times_S Y_s \to X \times_S Y)^*K \cong 
\Delta_{Y_s/k, *} \mathcal{O}_{Y_s}
$$
in $D(\mathcal{O}_{Y_s \times Y_s})$. Then there is an open neighbourhood
$s \in U \subset S$ such that $Y|_U$ is isomorphic to $Y_s \times U$ over $U$.
\end{lemma}

\begin{proof}
Denote $i : Y_s \times Y_s = X_s \times Y_s \to X \times_S Y$
the natural closed immersion. (We will write $Y_s$ and not $X_s$
for the fibre of $X$ over $s$ from now on.) Let
$z \in Y_s \times Y_s = (X \times_S Y)_s \subset X \times_S Y$
be a closed point. As indicated we think of $z$ both as a closed point
of $Y_s \times Y_s$ as well as a closed point of $X \times_S Y$.

\medskip\noindent
Case I: $z \not \in \Delta_{Y_s/k}(Y_s)$. Denote $\mathcal{O}_z$
the coherent $\mathcal{O}_{Y_s \times Y_s}$-module supported at $z$
whose value is $\kappa(z)$. Then $i_*\mathcal{O}_z$ is the
coherent $\mathcal{O}_{X \times_S Y}$-module supported at $z$
whose value is $\kappa(z)$. Our assumption means that
$$
K \otimes_{\mathcal{O}_{X \times_S Y}}^\mathbf{L} i_*\mathcal{O}_z =
Li^*K \otimes_{\mathcal{O}_{Y_s \times Y_s}}^\mathbf{L} \mathcal{O}_z = 0
$$
Hence by Lemma \ref{lemma-orthogonal-point-sheaf}
we find an open neighbourhood $U(z) \subset X \times_S Y$ of $z$
such that $K|_{U(z)} = 0$. In this case we set $Z(z) = \emptyset$
as closed subscheme of $U(z)$.

\medskip\noindent
Case II: $z \in \Delta_{Y_s/k}(Y_s)$. Since $Y_s$ is smooth over $k$
we know that $\Delta_{Y_s/k} : Y_s \to Y_s \times Y_s$ is a
regular immersion, see More on Morphisms, Lemma
\ref{more-morphisms-lemma-smooth-diagonal-perfect}.
Choose a regular sequence $\overline{f}_1, \ldots, \overline{f}_r \in
\mathcal{O}_{Y_s \times Y_s, z}$ cutting out the ideal sheaf of
$\Delta_{Y_s/k}(Y_s)$. Since a regular sequence is Koszul-regular
(More on Algebra, Lemma \ref{more-algebra-lemma-regular-koszul-regular})
our assumption means that
$$
K_z \otimes_{\mathcal{O}_{X \times_S Y, z}}^\mathbf{L}
\mathcal{O}_{Y_s \times Y_s, z}
\in D(\mathcal{O}_{Y_s \times Y_s, z})
$$
is represented by the Koszul complex on
$\overline{f}_1, \ldots, \overline{f}_r$ over
$\mathcal{O}_{Y_s \times Y_s, z}$.
By Lemma \ref{lemma-deform-koszul} applied to
$\mathcal{O}_{S, s} \to \mathcal{O}_{X \times_S Y, z}$
we conclude that $K_z \in D(\mathcal{O}_{X \times_S Y, z})$ is
represented by the Koszul complex on a regular sequence
$f_1, \ldots, f_r \in \mathcal{O}_{X \times_S Y, z}$
lifting the regular sequence
$\overline{f}_1, \ldots, \overline{f}_r$
such that moreover $\mathcal{O}_{X \times_S Y}/(f_1, \ldots, f_r)$
is flat over $\mathcal{O}_{S, s}$.
By some limit arguments (Lemma \ref{lemma-limit-arguments})
we conclude that there exists an affine open neighbourhood
$U(z) \subset X \times_S Y$ of $z$ and a closed subscheme
$Z(z) \subset U(z)$ such that
\begin{enumerate}
\item $Z(z) \to U(z)$ is a regular closed immersion,
\item $K|_{U(z)}$ is quasi-isomorphic to $\mathcal{O}_{Z(z)}$,
\item $Z(z) \to S$ is flat,
\item $Z(z)_s = \Delta_{Y_s/k}(Y_s) \cap U(z)_s$
as closed subschemes of $U(z)_s$.
\end{enumerate}

\noindent
By property (2), for $z, z' \in Y_s \times Y_s$, we
find that $Z(z) \cap U(z') = Z(z') \cap U(z)$ as closed subschemes.
Hence we obtain an open neighbourhood
$$
U = \bigcup\nolimits_{z \in Y_s \times Y_s\text{ closed}} U(z)
$$
of $Y_s \times Y_s$ in $X \times_S Y$ and a closed subscheme $Z \subset U$
such that (1) $Z \to U$ is a regular closed immersion,
(2) $Z \to S$ is flat, and (3) $Z_s = \Delta_{Y_s/k}(Y_s)$.
Since $X \times_S Y \to S$ is proper, after replacing $S$
by an open neighbourhood of $s$ we may assume $U = X \times_S Y$.
Since the projections $Z_s \to Y_s$ and $Z_s \to X_s$
are isomorphisms, we conclude that after shrinking $S$
we may assume $Z \to Y$ and $Z \to X$ are isomorphisms, see
Lemma \ref{lemma-isomorphism-in-neighbourhood}.
This finishes the proof.
\end{proof}

\end{document}
